\FloatBarrier
\section{FT inner shield and flange}

The fifth beamline study renames innerShieldAndFlange to FTInnerShieldAndFlange and optimizes its
engineering tessellation. The same rename applies to ten source and generated runtime CAD variations and ten
GEMC2 SQLite CAD rows. The tungsten-alloy material, density, color, mother, position, passive rotation,
coordinate frame, and overall bounds remain unchanged. Custom CAD placements must use the new name to match
the STL filename.

The original has 34,448 triangles, two closed manifold surface components, aggregate genus 17, and no
detected self-intersections or exact zero-area faces. The small inward-oriented component is an enclosed
pocket with signed volume -9.663980 mm$^3$. All 474 of its triangles are protected and verified in the saved
STL at exactly the original coordinates. Only the larger component is decimated, using topology-preserving
quadric collapse. No holes are filled, components removed, or surfaces reconstructed.

The saved output contains 22,000 triangles (36.14\% fewer) and 1,100,084 bytes instead of 1,722,484. It
retains both components, aggregate genus, and bounds; no self-intersections or zero-area triangles are
detected. Bidirectional vertex, edge, and surface sampling measures a maximum difference of 0.0481 mm, below
the 0.05 mm guard. Every vertex and approximately 200,000 edge/surface samples per direction are checked;
these sampled maxima are not proven global error bounds. Total signed volume changes by -0.00004391\%. A
separate component integration gives -0.00004114\% for the main component and no change for the pocket. The
builder reproduces the optimized STL byte for byte from the archived original.

Table~\ref{tab:inner-shield-meshes} summarizes the mesh inventory.

\begin{table}[htbp]
\centering
\caption{Original and optimized FT inner shield and flange.}
\label{tab:inner-shield-meshes}
\begin{tabular}{lrr}
\toprule
Quantity & Original & Optimized \\
\midrule
Triangles & 34,448 & 22,000 \\
Bytes & 1,722,484 & 1,100,084 \\
Closed surface components & 2 & 2 \\
Aggregate genus & 17 & 17 \\
Protected pocket faces & 474 & 474 \\
\bottomrule
\end{tabular}
\end{table}

The benchmark isolates this one CAD volume in air with the recorded GEMC3 executable and Geant4 stack,
FTFP\_BERT without optical processes, zero magnetic field, one thread, default production cuts, and no GUI
or event output. Both fixtures use the new volume name, identical physical placement attributes, and the
same alloy definition. Matched 5 GeV LUND primaries start at the target, uniform in solid angle over
0.25--3.0 degrees and full azimuth. Five alternating-order repetitions use seeds 12345--12349 and 200 events
per electron, proton, positive-pion, or neutron run. Separate zero-event runs measure startup without
overlap checking; time since beamOn measures transport. CPU time and peak RSS are retained separately.

Table~\ref{tab:inner-shield-timing} gives median transport wall seconds per 200 events.

\begin{table}[htbp]
\centering
\caption{Median transport wall seconds over five repetitions.}
\label{tab:inner-shield-timing}
\begin{tabular}{lrrr}
\toprule
Particle & Original & Optimized & Original/optimized \\
\midrule
e- & 18.565 & 15.048 & 1.234 \\
proton & 2.790 & 2.426 & 1.150 \\
pi+ & 3.655 & 3.292 & 1.110 \\
neutron & 2.504 & 2.425 & 1.033 \\
\bottomrule
\end{tabular}
\end{table}

Median startup simulation wall time changes from 0.9492 to 0.8435 s; peak loading RSS changes from 175.23 to
151.28 MiB. Median original/optimized transport wall-time ratios range from 1.033 to 1.234. All 52 processes
exit normally, and all forty transport runs report the 200 requested events. Results describe this isolated
workload and are not pooled with earlier beamline or LTCC benchmarks.

Table~\ref{tab:inner-shield-warnings} gives Geant4 exception counts over all 52 processes.

\begin{table}[htbp]
\centering
\caption{Geant4 exception counts for the isolated inner-shield samples.}
\label{tab:inner-shield-warnings}
\begin{tabular}{lrr}
\toprule
Diagnostic & Original & Optimized \\
\midrule
GeomSolids1001 & 26 & 0 \\
GeomNav1002 & 0 & 0 \\
GeomNav0003 & 0 & 0 \\
GeomVol1002 & 0 & 0 \\
\bottomrule
\end{tabular}
\end{table}

The diagnostic table records the original wrong-orientation import warning and all other Geant4 exceptions
in the sampled runs. Full default beamline construction completes without a missing mesh, with 4 solid
warnings and 11 overlap reports, compared with five and ten before this change. The full check reports no
overlap involving the optimized inner shield; the listed overlap reports concern other volumes. Their
sampled count can change when CAD loading order and surface tessellation change. All fifteen native
geometry/material variations and fourteen CAD variations pass the semantic comparison against GEMC2. Matched
primaries and seeds do not ensure identical stochastic shower histories, and completion does not establish
detector-response equivalence or complete overlap freedom.

\FloatBarrier
Figure~\ref{fig:inner-shield-stl} compares the stored triangulations.

\begin{figure}[htbp]
\centering
\includegraphics[width=\textwidth]{figures/inner_shield_stl_comparison.pdf}
\caption{Original (left) and optimized (right) FT inner shield and flange STL surfaces.
Matched cameras and bounds use original STL coordinates; the small enclosed pocket is retained exactly.}
\label{fig:inner-shield-stl}
\end{figure}
